\documentclass[10pt,a4paper]{amsart}
\usepackage[margin=19mm]{geometry}
\usepackage{amsmath,amssymb,mathtools}
\usepackage[scr=rsfs,cal=euler]{mathalfa}
\usepackage{xfrac}
\usepackage[unicode,hidelinks,bookmarks=false]{hyperref}
\newcommand{\ie}{i.e.\ }
\newcommand{\cf}{cf.\ }
\newcommand{\resp}{respectively\ }
\newcommand{\Subsection}{\subsection}
\providecommand{\nobreakdash}{\nobreak\hbox{-}}
\setlength{\emergencystretch}{2em}
\multlinegap0pt
\usepackage{bm}
\usepackage{bbm}
\usepackage{dsfont}
\usepackage{xspace}
\usepackage{cancel}
\usepackage{mathtools}
\usepackage{amsthm}
\makeatletter
\@ifundefined{theorem}{\newtheorem{theorem}{Theorem}}{}
\@ifundefined{proposition}{\newtheorem{proposition}[theorem]{Proposition}}{}
\makeatother
\def\FF{\mathbb{F}}
\def\ZZ{\mathbb{Z}}
\def\ini{\operatorname{in}}
\title[Generalized binomial edge ideals are Cartwright-Sturmfels]{Generalized binomial edge ideals are Cartwright-Sturmfels}
\author{Aldo Conca, Emanuela De Negri, Volkmar Welker}
\date{Source: arXiv:2512.22012v2 (CC BY 4.0)}
\begin{document}
\maketitle

\noindent\emph{Source and licence.} Generalized binomial edge ideals are Cartwright-Sturmfels. Version arXiv:2512.22012v2 (CC BY 4.0), distributed under the Creative Commons Attribution 4.0 International licence, as recorded in the arXiv licence metadata for this exact version and shown on the abstract page https://arxiv.org/abs/2512.22012v2. Full rights notice: licenses/arxiv-2512.22012-CC-BY-4.0.txt.\par\smallskip

\noindent\textit{Editorial scope.} Counted unit: the Theorem whose source label is \texttt{main} in the article (source0.tex lines 501--509, source label \texttt{main}), delivered together with the article's own complete original proof of it (source0.tex lines 510--545, one \texttt{proof} environment of 36 lines). No mathematics is written, repaired or re-derived: statement and proof are the article's own text line for line. Required context retained uncounted: easy2 (lines 385--392); gin (lines 473--476); cond (lines 480--483). Every definition, display and label of those lines is retained unchanged. Omitted from this package and not redistributed: 1--384, 393--472, 477--479, 484--500, 546--1139. Nothing inside a retained excerpt is omitted or altered: every retained body is delivered exactly as the article's lines are written, so no omission receipt of any kind accompanies this package. Citations: the retained excerpts carry no citation command at all, so no bibliography entry of the article is transcribed and no reference is redistributed. Reference adaptation: none was needed and none was made; every reference inside the delivered unit resolves inside it, so no unresolved reference or citation is produced. Printed slips kept literal: no printed word, symbol, hypothesis or number of the counted statement or of its proof was altered. Layout and class: the article's own class is \texttt{amsart} with packages including mathptmx, amssymb, color, xcolor, hyperref, tikz, doi; the wrapper is the lane's pinned \texttt{amsart} editorial wrapper at 10pt on A4, which restates the theorem-like environments the retained text needs and adds the article's own macro definitions that the retained text needs (\texttt{\textbackslash FF}, \texttt{\textbackslash ZZ}, \texttt{\textbackslash ini}), with extra wrapper packages (bm, bbm, dsfont, xspace, cancel, mathtools). The delivered numbering is the unit's own. Assets: no retained span contains a figure environment, a graphics-inclusion command, a PDF-page inclusion, a table or a TikZ picture; no image file is copied into this package, the package declares no asset dependency and no image file is redistributed with it. Licence: version arXiv:2512.22012v2 of this article is distributed under the Creative Commons Attribution 4.0 International licence, as recorded in the arXiv licence metadata for this exact version and shown on the abstract page https://arxiv.org/abs/2512.22012v2. A full rights notice is delivered at licenses/arxiv-2512.22012-CC-BY-4.0.txt. Characters: every retained excerpt is pure ASCII, so the delivered engine, XeLaTeX with its default Latin Modern text font, reports no missing character.
\begin{proposition}
\label{easy2} Let $\mathcal{F}\subseteq 2^{[n]}$ such that:
\begin{align} \label{eq:union}
    \text{ if } A,B\in \mathcal{F}  \text{ and } A\cap B\neq \emptyset \text{ then } A\cup B\in \mathcal{F}. 
    \end{align}
    Let $I$ be a CS ideal and  let $J$ be the ideal generated by the vector spaces $I_A$ with $A\in \mathcal{F}$. Then $J$ is CS. More precisely, for every term order $>$ the initial ideal  
$\ini_{>}(J)$ is generated  by the monomial vector spaces $\ini_{>}(I_A)$ with $A\in \mathcal{F}$. 
\end{proposition}

\begin{equation} 
\label{gin}
 \prod_{j\in A} x_{i_jj} 
 \end{equation}  

 \begin{equation} 
\label{cond} 
   i_j\in [m]   \mbox{ for every } j\in A \mbox{  and  } \sum_{j\in A} i_j\leq m(|A|-1).
  \end{equation}   

 \begin{theorem} 
 \label{main} Let $G = ([n],E)$ be a graph. 
 Then the generalized binomial edge ideal $I_G(m)$ is CS for all $m$. 
 More precisely, let $>$ be a term order satisfying
 $$x_{1j}>x_{2j} > \cdots > x_{mj} \text{ for } j \in [n].$$
 
 Then  the $\ZZ^n$-graded  generic initial ideal of $I_G(m)$ is generated by the monomials \eqref{gin} for the subsets $A \subseteq[n]$
 such that $G_A$ is connected and \eqref{cond} is satisfied.
\end{theorem}

\begin{proof} 
Set   $\mathcal F=\big\{\, A\subseteq [n]\, :\, G_A \text{ is connected }\,\}$ and  $I=I_2(X)$.  
If $A,B \in \mathcal F$ are
such that $A \cap B \neq \emptyset$ then $G_A$ and $G_B$ are two
connected graphs sharing at least one vertex. Thus 
$G_A \cup G_B = (A \cup B, E_A \cup E_B)$ is connected. Since
$E_A \cup E_B \subseteq E_{A \cup B}$ it follows that
$G_{A \cup B} = (A \cup B,E_{A\cup B})$ is connected and
$A \cup B \in \mathcal F$. In particular, $\mathcal F$
satisfies the assumptions of Proposition \ref{easy2}.

Let $J$ be the ideal generated by the vector spaces $I_A$ with $A\in \mathcal F$. By Proposition \ref{easy2} we know that $J$ is CS and 
its generic initial ideal is generated by the monomials \eqref{gin}
which satisfy \eqref{cond}. Therefore, it is enough to prove that $J=I_G(m)$. 

For the inclusion $J\supseteq I_G(m)$ observe that every edge $\{j,k\}$ of $G$ we have  $\{j,k\}\in \mathcal F$ and hence the generators of $I_G(m)$ are contained in $J$.  

For the other inclusion, we first observe the following. 
Denote by $C_{jk}$ the vector space  generated by the $2$-minors of $X$ using only columns $j$ and $k$ and by $C_\ell$ the vector space generated by the entries of column $\ell$. By double Laplace expansion one sees immediately that 

\begin{equation}
\label{expansion} 
 C_{jk} C_{\ell} \subseteq C_{\ell k} C_{j}  +  C_{j \ell} C_{k} .
\end{equation} 

Now let $A\in \FF$. We have to show that $I_A\subset I_G(m)$.   Since  $I_A$ is the sum of the vector spaces 
$$W_{j,k}:=C_{jk}\prod_{\ell\in A\setminus\{j,k\} }C_\ell$$ 
where $ \{j,k\}\subset A$, it is enough to show that $W_{j,k}\subseteq I_G(m)$.  As $G_A$ is connected, there exists a path $j=p_0,p_1,\dots, p_s=k$ in $G$ with $p_0,\ldots,p_s\in A$. 
We prove $W_{j,k} \subseteq I_G(m)$ by induction on $s$. 
If $s=1$ then $C_{jk} \subset I_G(m)$ and so $W_{jk} \subset  I_G(m)$. If $s>1$ then
by \eqref{expansion} we have
$C_{jk}C_{p_1}\subseteq C_{jp_1}C_{k}+C_{p_1k}C_{a}$. 
Thus  
$W_{jk} \subset W_{j,p_1}+W_{p_1,k}$. By induction on $s$, we may assume that both $W_{j,p_1}$ and $W_{p_1,k}$ are contained in $I_G(m)$ so that we can conclude that 
$W_{j,k}\subset   I_G(m)$. 
  \end{proof}  

\end{document}
